\subsection{Formal-codegree trace ideals}
\label{subsec:formal-codegree-trace-ideals}

The field-containment lemma admits a useful consequence even when the signed
dimension field has not been identified.  For an algebraic integer $f$, write
\[
\mathfrak t(f)=\operatorname{Tr}_{\mathbb Q(f)/\mathbb Q}
       \bigl(f\mathcal O_{\mathbb Q(f)}\bigr)=\tau(f)\mathbb Z,
\qquad \tau(f)>0.
\]
Conjugating $f$ does not change this integer ideal.

\begin{theorem}[Formal-codegree trace-ideal obstruction]
\label{thm:formal-codegree-trace-ideal}
Let $R$ be a commutative self-dual rank-six fusion ring, with regular matrices
$N_i$ and Casimir trace $T_0=\operatorname{Tr}(\sum_iN_i^2)$.
Choose any collection of its formal codegrees $f$ and put
$L=\operatorname{lcm}_f\tau(f)$.
If $R$ is categorifiable, the signed sphericalization has a dimension field
$K$ and global dimension $\Delta$ satisfying
\[
0<\operatorname{Tr}_{K/\mathbb Q}(\Delta)\in L\mathbb Z,
\qquad
\operatorname{Tr}_{K/\mathbb Q}(\Delta)\le T_0.
\]
In particular, $L>T_0$ excludes every complex fusion categorification.
No pivotality assumption or local-field isolation is required.
\end{theorem}
\begin{proof}
Lemma~\ref{lem:dimension-field-containment} gives $F=\mathbb Q(f)\subseteq K$
and $\Delta/f\in\mathcal O_K$.  The relative trace of an algebraic integer
is integral.  Therefore
\[
\operatorname{Tr}_{K/F}(\Delta)
=f\operatorname{Tr}_{K/F}(\Delta/f)\in f\mathcal O_F.
\]
Transitivity of trace gives the asserted divisibility for every selected $f$.
The embeddings of $K$ give distinct characters of the signed algebra, with
formal codegrees the corresponding conjugates of $\Delta$.  All its formal
codegrees are positive.  Thus their total, the signed Casimir trace, is at
least $\operatorname{Tr}_{K/\mathbb Q}(\Delta)$.  The signed coefficient
bounds give
$\operatorname{Tr}(\sum_iM_i^2)=\sum_{i,j,k}(t_{ij}^k)^2
\le\sum_{i,j,k}(N_{ij}^k)^2=T_0$.
\end{proof}

The implementation uses every rational and quadratic original codegree.
For a rational integer $f$, $\tau(f)=f$.  For a quadratic codegree
$f=(A+b\sqrt q)/2$, with $q>1$ squarefree, the integral basis gives
\[
\tau(f)=
\begin{cases}
\gcd\bigl(A,(A+bq)/2\bigr),&q\equiv1\pmod4,\\
\gcd(A,bq),&q\equiv2,3\pmod4.
\end{cases}
\]
Omitting higher-degree factors weakens the test and creates no exclusion.
The calculation depends only on the complete multiplication tensor.

For \texttt{R6\_8\_6997}, the rational formal codegrees include $15$, while
the FP codegree has minimal polynomial $X^2-6145X+30725$ and trace ideal
$6145\mathbb Z$.  Hence $L=18435>T_0=6173$, giving a general exclusion.
The same test supplies independent exclusions of
\texttt{R6\_7\_3812}, \texttt{R6\_8\_6234}, \texttt{R6\_8\_6509},
\texttt{R6\_8\_6940}, and \texttt{R6\_8\_6973}.

\begin{corollary}[Signed saturation from a trace ideal]
\label{cor:trace-ideal-saturation}
For an unordered triple of nonunit basis labels with coefficient $a\ge2$,
let $w$ be its number of ordered permutations, and put
\[
\eta=\min w\bigl(a^2-(a-2)^2\bigr),
\]
where the minimum runs over those triples.  If $0\le T_0-L<\eta$, every
signed sphericalization coefficient has absolute value equal to its
original nonnegative coefficient.  All remaining choices are finite sign
choices subject to associativity.
\end{corollary}
\begin{proof}
Parity forces any decrease of absolute value to be at least two.  Full
symmetry forces the same decrease on all $w$ ordered permutations.  Such a
decrease loses at least $\eta$ from the Casimir trace, whereas
Theorem~\ref{thm:formal-codegree-trace-ideal} permits loss at most $T_0-L$.
The unit coefficients do not change in the signed quotient.
\end{proof}

For \texttt{R6\_8\_6520}, one has $L=6145$, $T_0=6165$, and $\eta=24$.
Thus its signed quotient is saturated.  Exhaustive exact signed associativity
leaves one table up to the certified sign gauges; that table is gauge-equivalent
to the original ring.  The non-FP character orbits have zero entries and cannot
be dimension characters.  Consequently every categorification would have a
pseudo-unitary Galois conjugate.  Proposition~\ref{prop:quadratic-gauss-family}
excludes that possibility, completing the general noncategorifiability proof.

The authoritative implementation is \texttt{new\_trace\_ideal.py}.  It checks
all arithmetic from the tensor, replays signed enumeration when used, and
rejects altered certificates and omitted signed tables.  Its control report
uses the frozen initial forty residual tensors and all fifty-four established
positive models; the latter are all admitted by this necessary test.
